% Options for packages loaded elsewhere
\PassOptionsToPackage{unicode}{hyperref}
\PassOptionsToPackage{hyphens}{url}
\documentclass[
]{article}
\usepackage{xcolor}
\usepackage{amsmath,amssymb}
\setcounter{secnumdepth}{-\maxdimen} % remove section numbering
\usepackage{iftex}
\ifPDFTeX
  \usepackage[T1]{fontenc}
  \usepackage[utf8]{inputenc}
  \usepackage{textcomp} % provide euro and other symbols
\else % if luatex or xetex
  \usepackage{unicode-math} % this also loads fontspec
  \defaultfontfeatures{Scale=MatchLowercase}
  \defaultfontfeatures[\rmfamily]{Ligatures=TeX,Scale=1}
\fi
\usepackage{lmodern}
\usepackage[left=0.5in,right=0.5in,top=0.3in,bottom=0.3in]{geometry} % Reduced margins
\usepackage{ragged2e} % For justify environment
\usepackage{tikz} % For custom tags

% Define custom commands for CV tags and sections
\newcommand{\cvsection}[1]{\textbf{#1}}
\newcommand{\cvtag}[1]{\tikz[baseline=(char.base)]{\node[draw,rounded corners=3pt,inner sep=2pt,fill=gray!10] (char) {#1};}}
\newcommand{\divider}{\vspace{0.3em}\hrule\vspace{0.3em}}
\ifPDFTeX\else
  % xetex/luatex font selection
\fi
% Use upquote if available, for straight quotes in verbatim environments
\IfFileExists{upquote.sty}{\usepackage{upquote}}{}
\IfFileExists{microtype.sty}{% use microtype if available
  \usepackage[]{microtype}
  \UseMicrotypeSet[protrusion]{basicmath} % disable protrusion for tt fonts
}{}
\makeatletter
\@ifundefined{KOMAClassName}{% if non-KOMA class
  \IfFileExists{parskip.sty}{%
    \usepackage{parskip}
  }{% else
    \setlength{\parindent}{0pt}
    \setlength{\parskip}{6pt plus 2pt minus 1pt}}
}{% if KOMA class
  \KOMAoptions{parskip=half}}
\makeatother
\usepackage{longtable,booktabs,array}
\newcounter{none} % for unnumbered tables
\usepackage{calc} % for calculating minipage widths
% Correct order of tables after \paragraph or \subparagraph
\usepackage{etoolbox}
\makeatletter
\patchcmd\longtable{\par}{\if@noskipsec\mbox{}\fi\par}{}{}
\makeatother
% Allow footnotes in longtable head/foot
\IfFileExists{footnotehyper.sty}{\usepackage{footnotehyper}}{\usepackage{footnote}}
\makesavenoteenv{longtable}
\ifLuaTeX
  \usepackage{luacolor}
  \usepackage[soul]{lua-ul}
\else
  \usepackage{soul}
\fi
\setlength{\emergencystretch}{3em} % prevent overfull lines
\providecommand{\tightlist}{%
  \setlength{\itemsep}{0pt}\setlength{\parskip}{0pt}}
\usepackage{bookmark}
\IfFileExists{xurl.sty}{\usepackage{xurl}}{} % add URL line breaks if available
\urlstyle{same}
\hypersetup{
  hidelinks,
  pdfcreator={LaTeX via pandoc}}

\author{}
\date{}

\begin{document}

\begin{center}{\Large\textbf{Hemanth Kumar J N}}\end{center}

\begin{center}
Bengaluru 560010, India \textbar{} hemanthgunduraj77@gmail.com
\textbar{} +91-6361613890 \textbar{}
\href{https://github.com/hemgowda}{\ul{GitHub}} \textbar{}
\href{https://linkedin.com/in/Hemanthakumar}{\ul{LinkedIn}}
\end{center}

\begin{center}\textbf{Summary}\end{center}
\vspace{-0.5em}

Backend \& Data Engineer with 4+ years of experience building LLM-powered
data platforms, ETL pipelines, and multi-cloud infrastructure across AWS
and GCP in the healthcare domain. Skilled in Snowflake, BigQuery, Python
backend services, and ML/GenAI pipelines. M.S. in Data Science with hands-on expertise in DevOps automation, CI/CD, and scalable
cloud-native architectures.
\vspace{-1.2em}

\begin{center}\cvsection{Skills}\end{center}
\vspace{-0.5em}

\begin{minipage}[t]{0.5\textwidth}
\textbf{Core Expertise}
\vspace{0.21em}

\textbf{Programming : }\cvtag{Python} \cvtag{SQL} \cvtag{PySpark} \cvtag{Data-Structures}
\vspace{0.09em}

\textbf{Backend \& Systems:}\cvtag{REST APIs} \cvtag{Distributed Systems}           \cvtag{Microservices}  \cvtag{System Design}  \cvtag{Event - Driven Architecture} 
\vspace{0.1em}

\textbf{Cloud Engineering :} \cvtag{AWS} \cvtag{GCP} \cvtag{Snowflake}  \cvtag{ETL/ELT} 
\vspace{0.1em}

\textbf{GenAI \& ML:}\cvtag{LLM} \cvtag{RAG } \cvtag{Vertex AI} \cvtag{Model Deployment } 

\end{minipage}
\hfill
\begin{minipage}[t]{0.48\textwidth}
\textbf{Tools \& Technologies}
\vspace{0.2em}

\textbf{Databases :} \cvtag{MySQL} \cvtag{MongoDB} \cvtag{BigQuery} 
\cvtag{Teradata}
\vspace{0.2em}

\textbf{DevOps \& Streaming:} \cvtag{Git} \cvtag{Docker} \cvtag{K8S} \cvtag{CI/CD} \cvtag{Jenkins} \cvtag{Terraform} \cvtag{CloudFormation} \cvtag{SQS} \cvtag{SNS} \cvtag{Linux}\vspace{0.2em}

\textbf{Data Science \& Analytics:} \cvtag{Scalable ML Pipelines} \cvtag{Data Modeling} \cvtag{Feature Engineering} \cvtag{Statistical Analysis} 
\end{minipage}
\vspace{-0.1em}

\begin{center}\textbf{Experience}\end{center}
\vspace{-0.7em}

\textbf{Carelon Global Solutions Private Limited} \hfill Jun 2022 -- Present

\emph{\textbf{Software Engineer III -- Cloud Data Engineering}}
\emph{Bengaluru, India}

\begin{itemize}
\item
  Architected Python backend pipelines integrating RDS and Snowflake
  metadata for job health analytics, and integrated LLM endpoints with
  retrieval and embedding workflows to power a natural-language chatbot
  -- reducing ad-hoc analyst requests by \textasciitilde60\%.
\item
  Built an ETL pipelines using SQL and Python with
  partitioning and clustering in Snowflake and RDS, improving query
  performance by 30--40\%; deployed Flask-based REST APIs on AWS EC2
  with auto-scaling achieving 99.9\% uptime.
\item
  Engineered a YAML-based automation framework cutting 80\% of manual
  operational work, and built a Python SQL-generation tool (COT, DDL,
  Falcon SQL) enabling creation of 1,000+ tables across Snowflake --
  recognized with Carelon PLT Appreciation Award.
\item
  Automated CI/CD deployments via AWS pipelines with CloudWatch
  monitoring; Dockerized services for scalable deployments and enforced
  data governance through centralized health and usage metadata
  tracking.
\item
  Designed and exposed LLM-based monitoring dashboards via Flask REST
  APIs, integrating governance rules and job-failure alerts with
  BigQuery and Snowflake metadata -- enabling stakeholders to query
  pipeline status in real time.
\end{itemize}
\vspace{-0.2em}

\begin{center}\textbf{Projects}\end{center}
\vspace{-0.5em}

\textbf{Data Health LLM Chatbot \& Analytics Platform} \hfill \emph{jun 2022 -- Present}
\vspace{-0.5em}

\emph{Python \textbar{} Pandas \textbar{} RDS \textbar{} BigQuery
\textbar{} Snowflake \textbar{} SQL \textbar{} LLM \textbar{} Embeddings
\textbar{} Flask}

\begin{itemize}
\item
  Built Python backend pipelines extracting metadata from RDS and
  BigQuery into consolidated monitoring tables; integrated LLM-based
  chatbot with retrieval and embedding workflows to answer natural
  language queries, cutting manual monitoring effort by
  \textasciitilde60\%.
\item
  Centralized data governance and table health/usage insights, enabling
  proactive alerting on pipeline failures and improving platform
  reliability.
\end{itemize}

%\textbf{ML \& Data Engineering with Vertex AI -- Healthcare Fraud Detection} \hfill %\emph{Mar 2022 -- Present}
%
%\emph{GCP \textbar{} Vertex AI \textbar{} BigQuery \textbar{} Cloud
%Storage \textbar{} Dataflow \textbar{} Pub/Sub \textbar{} Looker Studio
%\textbar{} XGBoost \textbar{} AutoML}
%
%\begin{itemize}
%\item
%  Designed an end-to-end 3-layer ML pipeline (Raw → Curated BigQuery →
%  Vertex AI Serving) for Medicaid/Medicare fraud detection using
%  XGBoost/AutoML; ingested streaming data via Pub/Sub and Dataflow
%  (Apache Beam).
%\item
%  Delivered real-time Looker Studio dashboards for claims monitoring and
%  enforced HIPAA compliance via IAM roles and VPC Service Controls.
%\end{itemize}

\textbf{Multi-Cloud Infrastructure Automation (AWS + GCP)} \hfill \emph{Mar 2024 -- Present}
\vspace{-0.5em}

\emph{AWS: EC2 \textbar{} VPC \textbar{} ELB \textbar{} WAF \textbar{}
RDS \textbar{} CloudFormation \textbar{} GCP: Compute Engine \textbar{}
Cloud SQL \textbar{} Terraform}

\begin{itemize}
\item
  Designed a secure multi-cloud architecture (AWS VPC + GCP VPC peering)
  with auto-scaling via EC2 Auto Scaling + GCP Instance Groups,
  CloudFront + Cloud CDN for global caching, and infrastructure-as-code
  using CloudFormation and Terraform for reproducible deployments.
\item
  Enhanced security posture by implementing AWS WAF and GCP Cloud Armor
  for application-level threat protection; automated provisioning
  ensured consistent, auditable environment configurations across dev,
  staging, and production.
\end{itemize}
\vspace{-0.3em}

\begin{center}\textbf{Education}\end{center}
\vspace{-0.3em}

\textbf{Manipal Academy of Higher Education, Bangalore} 2023 -- 2025

\emph{M.S. in Data Science -- SGPA: 9.4}

%\textbf{Jawaharlal Nehru National College of Engineering, Shivamogga}
%2018 -- 2022
%
%\emph{B.E. in Information Science and Engineering -- CGPA: 8.1}
\vspace{-0.9em}

\begin{center}\cvsection{Certifications}\end{center}
\vspace{-0.5em}

\begin{minipage}[t]{0.48\textwidth}
\textbf{Google Certified Associate Cloud Engineer} -- Google
\vspace{0.1em}

\textbf{Machine Learning \& Data Science Capstone} -- IBM
\vspace{0.1em}

\textbf{Databases \& SQL for Data Science with Python} -- IBM
\end{minipage}
\hfill
\begin{minipage}[t]{0.48\textwidth}
\textbf{Building Statistical Models in R} -- Coursera
\vspace{0.1em}

\textbf{Object-Oriented Programming in Python} -- Infytq
\vspace{0.1em}

\textbf{Programming Fundamentals using Python} -- Infytq
\end{minipage}

\end{document}
